\documentclass[12pt,a4paper]{article}
\usepackage[utf8]{inputenc}
\usepackage{amsmath}
\usepackage{amssymb}
\usepackage{geometry}
\usepackage{listings}
\usepackage{xcolor}
\usepackage{hyperref}
\usepackage{enumitem}
\usepackage{booktabs}
\usepackage{graphicx}

\geometry{margin=1in}

\title{FedAvg (Federated Averaging): Giải thích chi tiết kèm công thức\\Hệ thống Federated Learning phát hiện DDoS}
\author{}
\date{\today}

\lstset{
    basicstyle=\ttfamily\small,
    keywordstyle=\color{blue},
    stringstyle=\color{red},
    commentstyle=\color{green!50!black},
    breaklines=true,
    frame=single
}

\begin{document}

\maketitle

\section{Tổng quan}

\textbf{Federated Averaging (FedAvg)} là thuật toán tối ưu hoá kinh điển được sử dụng trong Federated Learning (FL).
Thuật toán này huấn luyện một \textbf{mô hình toàn cục (global model)} trên nhiều client/worker mà không cần chuyển dữ liệu thô về server.
Server điều phối quá trình huấn luyện theo \textbf{vòng (round)}. Ở mỗi vòng:
\begin{enumerate}
    \item Server phát các tham số toàn cục hiện tại tới các client được chọn.
    \item Mỗi client thực hiện \textbf{huấn luyện cục bộ (local training)} (thường dùng SGD/Adam trên dữ liệu riêng).
    \item Client gửi tham số đã cập nhật về server.
    \item Server tổng hợp các cập nhật bằng \textbf{trung bình có trọng số}.
\end{enumerate}

Trong repo của bạn, Flower sử dụng strategy \texttt{FedAvg} (xem \texttt{flower\_server.py}), và server trong \texttt{flower\_server\_metrics.py} kế thừa nó thành \texttt{MetricsFedAvg} để theo dõi metric cho dashboard/MLflow.

\section{Ký hiệu}

\begin{itemize}
    \item $K$: tổng số client/worker.
    \item $k \in \{1,\dots,K\}$: chỉ số client.
    \item $n_k$: số mẫu huấn luyện trên client $k$.
    \item $n = \sum_{k=1}^{K} n_k$: tổng số mẫu huấn luyện trên tất cả client.
    \item $w \in \mathbb{R}^d$: tham số mô hình (trọng số).
    \item $t \in \{0,1,2,\dots\}$: chỉ số vòng FL.
    \item $S_t \subseteq \{1,\dots,K\}$: tập client tham gia ở vòng $t$ (tham gia một phần).
    \item $\ell(w; x, y)$: hàm mất mát cho từng mẫu (ví dụ: cross-entropy).
    \item $\eta$: learning rate cục bộ.
    \item $E$: số epoch huấn luyện cục bộ mỗi vòng.
    \item $B$: kích thước mini-batch cục bộ.
\end{itemize}

\section{Mục tiêu tối ưu}

FedAvg hướng tới việc tối thiểu hoá \textbf{rủi ro thực nghiệm toàn cục} (có trọng số theo kích thước dữ liệu từng client):
\[
\min_{w \in \mathbb{R}^d} F(w) \quad \text{where} \quad
F(w) = \sum_{k=1}^{K} \frac{n_k}{n}\,F_k(w).
\]

Mỗi client có một \textbf{hàm mục tiêu cục bộ}:
\[
F_k(w) = \frac{1}{n_k}\sum_{i=1}^{n_k} \ell\bigl(w; x_{k,i}, y_{k,i}\bigr).
\]

\textbf{Điểm quan trọng:} server không bao giờ nhìn thấy $(x_{k,i}, y_{k,i})$; server chỉ nhận cập nhật tham số mô hình.

\section{Thuật toán FedAvg (dạng toán học)}

\subsection{Cập nhật phía server (trung bình có trọng số)}

Ở vòng $t$, server giữ trọng số toàn cục $w_t$ và chọn tập client tham gia $S_t$.
Sau huấn luyện cục bộ, mỗi client gửi về mô hình đã cập nhật $w_{t+1}^{(k)}$.
Quy tắc tổng hợp FedAvg là:
\[
w_{t+1} = \sum_{k \in S_t} \frac{n_k}{\sum_{j \in S_t} n_j}\, w_{t+1}^{(k)}.
\]

Đây là \textbf{trung bình có trọng số theo số mẫu}. Client có nhiều dữ liệu sẽ ảnh hưởng tới mô hình toàn cục nhiều hơn.

\subsection{Huấn luyện cục bộ phía client (Local SGD / local optimizer)}

Mỗi client tham gia sẽ khởi tạo mô hình cục bộ bằng mô hình từ server:
\[
w_{t,0}^{(k)} \leftarrow w_t.
\]

Sau đó client huấn luyện $E$ epoch cục bộ. Với mini-batch ở bước cục bộ $s$:
\[
w_{t,s+1}^{(k)} = w_{t,s}^{(k)} - \eta \,\nabla \ell\left(w_{t,s}^{(k)}; \mathcal{B}_{k,s}\right),
\]
trong đó $\mathcal{B}_{k,s}$ là một mini-batch lấy mẫu từ dữ liệu của client $k$.

Sau khi hoàn tất huấn luyện cục bộ, client $k$ trả về:
\[
w_{t+1}^{(k)} \leftarrow w_{t,S_k}^{(k)},
\]
trong đó $S_k$ là tổng số bước cập nhật cục bộ trên client $k$.

\subsection{Góc nhìn khác: trung bình hoá cập nhật (Delta)}

Đôi khi, việc biểu diễn FedAvg theo độ lệch tham số (delta) sẽ tiện hơn:
\[
\Delta_t^{(k)} = w_{t+1}^{(k)} - w_t.
\]
Khi đó, cập nhật tổng hợp là:
\[
w_{t+1} = w_t + \sum_{k \in S_t} \frac{n_k}{\sum_{j \in S_t} n_j}\,\Delta_t^{(k)}.
\]

\section{Giả mã (dễ đọc)}

\begin{lstlisting}[language={},caption={Giả mã FedAvg (server + clients)}]
Server initializes global model weights w0
for round t = 0, 1, 2, ..., T-1 do:
    St <- sample_clients(K, fraction_fit)   # participating clients
    broadcast wt to all k in St

    # Local training happens in parallel on clients
    for each client k in St (in parallel) do:
        w <- wt
        for epoch e = 1..E do:
            for minibatch B in client_data(k) do:
                w <- w - eta * grad(loss(w; B))
        return wk <- w, and nk (# local samples)

    # Server aggregation (weighted average)
    wt+1 <- sum_{k in St} (nk / sum_{j in St} nj) * wk
end for
\end{lstlisting}

\section{Vì sao FedAvg hoạt động (trực giác)}

\subsection{Trường hợp IID (lý tưởng)}
Nếu dữ liệu trên mỗi client là IID từ cùng một phân phối, thì gradient cục bộ là ước lượng không chệch của gradient toàn cục:
\[
\mathbb{E}[\nabla F_k(w)] \approx \nabla F(w).
\]
Trong chế độ này, FedAvg xấp xỉ rất gần với mini-batch SGD tập trung, và hội tụ thường nhanh và ổn định.

\subsection{Trường hợp Non-IID (trôi client / client drift)}
Với dữ liệu Non-IID, gradient giữa các client có thể lệch hướng nhau:
\[
\nabla F_k(w) \not\approx \nabla F(w).
\]
Huấn luyện cục bộ nhiều epoch $E$ có thể gây \textbf{client drift} (mô hình cục bộ bị kéo về cực tiểu địa phương), khiến phép trung bình kém hiệu quả và đôi khi không ổn định.

Đây là lý do vì sao chia dữ liệu Non-IID (ví dụ chia lệch lớp trong \texttt{flower\_worker.py}) có thể làm giảm độ chính xác hoặc làm chậm hội tụ.

\section{Siêu tham số quan trọng và ảnh hưởng}

\begin{table}[h]
\centering
\begin{tabular}{@{}lll@{}}
\toprule
\textbf{Tham số} & \textbf{Ý nghĩa} & \textbf{Ảnh hưởng thường gặp} \\
\midrule
$T$ & số vòng FL & Nhiều vòng $\rightarrow$ hội tụ tốt hơn, tốn giao tiếp hơn \\
$E$ & số epoch cục bộ mỗi vòng & $E$ lớn giảm giao tiếp nhưng tăng drift (Non-IID) \\
$\eta$ & learning rate cục bộ & Quá lớn $\rightarrow$ phân kỳ; quá nhỏ $\rightarrow$ chậm \\
$B$ & batch size cục bộ & $B$ lớn giảm nhiễu gradient nhưng có thể làm update chậm hơn \\
fraction\_fit & tỷ lệ client tham gia & Tỷ lệ thấp giảm chi phí/vòng nhưng tăng phương sai \\
weights & trọng số tổng hợp & Trọng số theo $n_k$ tăng hiệu quả thống kê \\
\bottomrule
\end{tabular}
\caption{Các siêu tham số của FedAvg}
\end{table}

\section{Liên hệ với triển khai Flower trong repo}

\subsection{FedAvg trong Flower}

Server của bạn dùng strategy \texttt{FedAvg} có sẵn trong Flower, cấu hình với các tham số như:
\begin{itemize}
    \item \texttt{fraction\_fit}: tỷ lệ client được lấy mẫu mỗi vòng,
    \item \texttt{min\_fit\_clients}: số client tối thiểu cần có để huấn luyện,
    \item \texttt{min\_available\_clients}: số client tối thiểu phải ``online/available''.
\end{itemize}

File \texttt{flower\_server\_metrics.py} kế thừa \texttt{FedAvg} thành \texttt{MetricsFedAvg} để log:
thời gian mỗi vòng, bytes gửi/nhận, số tham số, và metric đánh giá đã tổng hợp.

\subsection{FedAvg tổng hợp ở đâu}

Trong Flower, phép tổng hợp được thực hiện trong hàm \texttt{aggregate\_fit} của strategy.
Class tuỳ biến của bạn gọi \texttt{super().aggregate\_fit(...)} (thực hiện FedAvg) và sau đó bổ sung phần logging/metrics.

\begin{lstlisting}[language=Python,caption={Flower strategy pattern used in your repo (simplified)}]
from flwr.server.strategy import FedAvg

class MetricsFedAvg(FedAvg):
    def aggregate_fit(self, rnd, results, failures):
        aggregated_weights, aggregated_metrics = super().aggregate_fit(rnd, results, failures)
        # ... your code: track bytes, params, iid flag, dashboard, mlflow ...
        return aggregated_weights, aggregated_metrics
\end{lstlisting}

\section{Lưu ý thực tế (triển khai thật)}

\subsection{Tham gia một phần (partial participation)}
Thực tế thường chỉ một phần client tham gia ở mỗi vòng. FedAvg hỗ trợ tự nhiên thông qua $S_t$.
Tuy nhiên việc lấy mẫu có thể làm tăng phương sai; trọng số theo $n_k$ giúp giảm nhưng không triệt tiêu hoàn toàn tính ngẫu nhiên.

\subsection{Client chậm (straggler) và rớt kết nối (dropout)}
Client có thể chậm hoặc không gửi kết quả về đúng hạn.
Thực tế server thường đặt timeout và tổng hợp các kết quả nhận được (Flower có hỗ trợ round timeout).

\subsection{Chi phí giao tiếp}
Chi phí giao tiếp mỗi vòng tăng theo kích thước mô hình $d$ và số client tham gia $|S_t|$.
Repo của bạn đã theo dõi:
\[
\text{Bytes per round} \approx \text{Bytes sent} + \text{Bytes received}.
\]
Giảm kích thước mô hình, nén trọng số, hoặc giảm tỷ lệ tham gia có thể giảm chi phí.

\subsection{Tăng cường riêng tư (tuỳ chọn)}
FedAvg tự thân không đảm bảo riêng tư tuyệt đối. Các kỹ thuật bổ trợ thường gặp gồm:
\begin{itemize}
    \item \textbf{Differential Privacy (DP)}: thêm nhiễu vào gradient/cập nhật hoặc dùng DP-SGD ở client.
    \item \textbf{Secure Aggregation}: tổng hợp bằng mật mã để server không thể xem cập nhật từng client.
\end{itemize}
Worker của bạn có tuỳ chọn DP; DP sẽ làm thay đổi động học tối ưu, nhưng cấu trúc ``trung bình ở server'' vẫn tương tự.

\section{Tóm tắt}

FedAvg giải bài toán:
\[
\min_w \sum_{k=1}^{K}\frac{n_k}{n}F_k(w)
\]
bằng cách lặp xen kẽ giữa:
\begin{enumerate}
    \item \textbf{Tối ưu cục bộ} ở client (nhiều bước SGD),
    \item \textbf{Trung bình có trọng số} ở server:
    \[
    w_{t+1} = \sum_{k \in S_t} \frac{n_k}{\sum_{j \in S_t} n_j} w_{t+1}^{(k)}.
    \]
\end{enumerate}

FedAvg hoạt động tốt nhất trong bối cảnh IID và có thể suy giảm khi Non-IID mạnh do hiện tượng client drift, điều này liên quan trực tiếp tới các thí nghiệm IID vs Non-IID của bạn.

\end{document}


